\chapter*{引言}
\addcontentsline{toc}{chapter}{引言}
\markboth{引言}{引言}
\origpage{4}{III}
本书旨在使非专门研究者更便于使用\term{$\ell$ 进上同调}{$\ell$-adic cohomology}{cohomologie $\ell$-adique}。我希望，读者借助本书，往往可以不必求助于 SGA 4 和 SGA 5 中那些枝蔓繁复的报告。书中也包含一些新结果。

第一篇由 J. F. Boutot 执笔，概览 SGA 4。它介绍主要结果及其证明思路；所取的一般性尽可能低，对于应用往往还不够。要了解完整的结果或详细的证明，SGA 4 仍然不可替代。

《迹公式报告》给出了 Frobenius 自身态射的\term{迹公式}{trace formula}{formule des traces}的完整证明。这就是 Grothendieck 在 SGA 5 中给出的证明，删去了所有不必要的细节。有了这篇报告，使用者应当可以把 SGA 5 放在一旁，将它视为一系列旁逸斜出的讨论，其中有些很有趣。也正因为有了这篇报告，SGA 5 即将可以按现状出版。《迹公式对三角和的应用》是它的补篇，说明如何运用迹公式研究\term{三角和}{trigonometric sum}{somme trigonométrique}，并给出若干例子。

其余各篇所面向的读者更为专门，行文风格也因此有所不同。《模 $\ell^n$ 与模 $p$ 的 $L$ 函数》以 SGA 4 XVII 5.5 对\term{对称幂}{symmetric power}{puissance symétrique}的研究为基础，对《报告》作“模意义下”的推广。《代数闭链所对应的上同调类》在不同背景下定义这一类，并说明相交与\term{杯积}{cup product}{cup-produit}之间的相容性。《对偶》汇集了几个已有但缺乏可引文献的结果，以及一些相容性关系。《$\ell$ 进上同调的有限性定理》是新的；其中尤其在无支撑限制的上同调中，给出了与紧支撑上同调中已知结果相应的有限性定理。

各篇的进一步说明，请见各自的引言。

\origpage{5}{IV}
最后，感谢 J. L. Verdier 允许我在此重印他的讲义《导出范畴（初稿）》。我认为，这份讲义至今仍很有用，而此前已很难找到。

本书内部引用各篇时，采用目录中以方括号标出的简称。

\vspace{8mm}
\begin{flushright}
Bures-sur-Yvette，1976 年 9 月 20 日\\[4mm]
Pierre Deligne
\end{flushright}

\clearpage
\section*{各篇的引用简称}
\begin{longtable}{@{}p{.23\linewidth}p{.70\linewidth}@{}}
\toprule
\textbf{简称} & \textbf{篇名}\\\midrule
{[Arcata]} & \hyperref[I]{平展上同调的出发点}（J. F. Boutot 执笔）\\[3pt]
{[Rapport]} & \hyperref[II]{迹公式报告}\\[3pt]
{[Fonction L mod.$\ell^n$]} & \hyperref[III]{模 $\ell^n$ 与模 $p$ 的 $L$ 函数}\\[3pt]
{[Cycle]} & \hyperref[IV]{代数闭链所对应的上同调类}（A. Grothendieck 著，P. Deligne 执笔）\\[3pt]
{[Dualité]} & \hyperref[V]{对偶}\\[3pt]
{[Sommes trig.]} & \hyperref[VI]{迹公式对三角和的应用}\\[3pt]
{[Th. finitude]} & \hyperref[VII]{$\ell$ 进上同调的有限性定理}（附 L. Illusie 撰写的附录）\\[3pt]
{[C.D.]} & \hyperref[VIII]{导出范畴（初稿）}（J. L. Verdier）\\
\bottomrule
\end{longtable}
